% Revisión de Fase 3, 31-10-2026.
\section{Requisitos no funcionales}
\label{sec:rnf}
Se conserva la identificación RNF de la matriz MoSCoW. Implementación no equivale
a cumplimiento medido: los estados siguientes se basan en código local y evidencia
histórica identificada, no en una auditoría actual de producción.
\begin{longtable}{p{1.6cm}p{2.1cm}p{10.5cm}}
\toprule ID & Estado & Fundamento y brecha \\\midrule\endhead
RNF-01 & Parcial & Interfaz de descubrimiento y cuenta implementada. Revalidar defectos de usabilidad del ciclo 11-10-2026. \\
RNF-02 & Parcial & Navegación adaptativa y proyecto Capacitor con APK debug existente. Faltan pruebas nativas en dispositivo y regresión responsive. \\
RNF-03 & Sin medición actual & El tiempo HTTP histórico no mide renderizado ni Core Web Vitals. Faltan volumen, concurrencia y métricas de interacción. \\
RNF-04 & Parcial & Auth, funciones Edge, RLS y secretos de servidor. Registro principal sin verificación de propiedad del buzón; revisar abuso y exposición de datos. La presencia de políticas no acredita su eficacia remota. \\
RNF-05 & Parcial & Guard y servidor exigen profiles.role=admin y proveedor email; visitor/editor/OAuth quedan fuera. Timeout de 10 min en cliente. Faltan pruebas negativas y de revocación. \\
RNF-06 & Parcial & Vercel y evidencia histórica de CI. No se aporta medición de disponibilidad actual ni SLA; el scheduler puede retrasarse. \\
RNF-07 & Diseño implementado & Paginación y upsert en componentes documentados. No se acredita capacidad bajo carga sin medición. \\
RNF-08 & Parcial & Pipeline de saneamiento y deduplicación documentado; revisar muestra actual del catálogo y ubicaciones aproximadas. \\
RNF-09 & Parcial & Servicios, pruebas Vitest, fuentes de diagramas y documentación versionable. Completar ejecución actual y trazabilidad de entrega. \\
RNF-10 & Disponible & Contenido en español y regiones/comunas chilenas; revisión editorial final pendiente. \\
RNF-11 & Parcial & Build web, configuración Vercel y Capacitor presentes; validar reproducción en entorno limpio. Docker no adoptado. \\
RNF-12 & Parcial & Roles ARIA, foco y movimiento reducido en formularios. No se declara conformidad WCAG 2.2 completa sin auditoría y recorrido manual. \\
\bottomrule
\caption{Estado de los RNF al corte documental.}\label{tab:rnfestado}
\end{longtable}

\subsection{Seguridad y limitaciones}
La autenticación comprueba identidad; la autorización administrativa se decide con
\texttt{profiles.role} y se repite en \texttt{admin-eventos}. Ocultar una ruta no
la protege. La clave pública del navegador no es un secreto y exige políticas RLS
correctas; \texttt{service\_role} no debe incorporarse al frontend, capturas ni Git.
El registro sin verificación de correo es una decisión funcional: confirmar mediante
Admin API no acredita que la persona controle el buzón. El correo de bienvenida no
reemplaza esa comprobación. El cierre por inactividad es cliente y puede reiniciarse
con recarga; no se describe como revocación de servidor.

Las notificaciones por correo tienen código en funciones, pero requieren configuración,
recepción comprobada y, para recomendaciones periódicas, invocación programada verificable.
La moderación de eventos está implementada; los reportes ciudadanos y la administración
de organizadores no se dan por completos por existir tablas o rutas. La ocultación
local de fuentes no cambia la ejecución del scraper. La estadística del dashboard es
un cálculo del catálogo, no analítica comercial.

\section{Pruebas y evidencia}
\label{sec:pruebas}
\subsection{Resultados históricos conservados}
El ciclo QA de Vercel del 11-10-2026 registra 13 escenarios: 6 aprobados,
6 con observaciones y 1 fallido, con 13 defectos abiertos. Se conserva en
\texttt{Evidencias\_QA\_Frontend\_2026-10-11} y en la matriz metodológica.
No se cierran defectos sin una nueva ejecución.
La evidencia de CI del scraper conserva cinco ejecuciones trazadas y el run 25
funcional con dos artefactos; su inicio fuera de ventana se registró como DEF-CI-001.
El inventario histórico describe 14 módulos de pruebas del scraper; no se certifica
que ese conteo represente la rama actual.

El anexo previo registra \texttt{ng test --watch=false} con 23/23 el 27-10-2026
y builds anteriores. Son antecedentes, no resultados reproducidos en esta revisión
documental. Los archivos actuales incluyen pruebas de recomendaciones, Supabase y
acciones de interés. No se infiere cobertura de seguridad a partir de su existencia.

\subsection{Plan incremental de cierre}
El Plan de pruebas metodológico incorpora casos F3-01 a F3-10 con resultados
\textbf{pendientes}. Incluyen registro/bienvenida, recuperación, perfil, roles,
OAuth, inactividad, moderación, persistencia, Android y regresión.
Los comandos previstos son \texttt{npm test -- --watch=false},
\texttt{npm run build} y las pruebas focalizadas de cada integración.
Una ejecución debe identificar versión, entorno, responsable y resultado real;
la compilación de esta documentación únicamente acredita generación de entregables.

\section{Despliegue sin Docker}
\label{sec:docker}
La estrategia adoptada no utiliza Docker: el scraper se ejecuta mediante GitHub
Actions, la SPA en Vercel y Android con Capacitor. Supabase es un servicio administrado.
No se agrega un despliegue alternativo ni se declara Docker como implementado.
Su adopción futura requeriría una decisión explícita de arquitectura y pruebas propias.
El Manual técnico describe la reproducción del build y las dependencias de servidor.

\section{Innovación y alcance del MVP}
\label{sec:innovacion}
El valor del proyecto es reunir fuentes culturales heterogéneas y conservar procedencia,
normalización, deduplicación y vigencia en un catálogo consultable. El uso selectivo de
NVIDIA NIM para OCR y enriquecimiento evita repetir trabajo sin cambios; los checkpoints
permiten reanudar procesos. Las recomendaciones al ciudadano se basan en afinidad y
señales explícitas; no se presentan como inferencia de un LLM.

El prototipo académico depende de servicios externos y cuotas configuradas; no dispone
de una infraestructura propia de inferencia. No se afirman ahorros, mejoras de latencia,
usuarios activos ni cifras de eventos actuales sin medición fechada.
El mapa diferencia coordenadas conocidas y ubicaciones aproximadas; esas aproximaciones
no deben presentarse como direcciones exactas.

\subsection{Sostenibilidad: propuesta futura}
El MVP es gratuito y no implementa pagos, suscripciones, venta de datos personales ni
comisiones por entradas. Un eventual contenido patrocinado exigiría validación legal,
comercial y ética, identificación visible, separación de recomendaciones orgánicas y
compatibilidad con las condiciones del hosting. No se certifican condiciones comerciales
vigentes de un proveedor mediante este documento. Cualquier monetización queda fuera
del alcance del cierre actual.
